\documentclass[12pt]{article}
\usepackage[utf8]{inputenc}
\usepackage[T1]{fontenc}
\usepackage{hyperref}
\title{Mesme Language Resources for K-12 Curriculum Development}
\author{Compiled by Claude Sonnet 5.5 (claude-sonnet-5-5) for Andr\'as Kornai}
\date{}
\begin{document}
\maketitle

\section{Introduction}
Mesme (ISO 639-3: zim; Glottocode: mesm1239), also called Zime or Djime, is a Chadic (Afro-Asiatic) language of southwestern Chad, in the Tandjil\'e region around K\'elo. Wikipedia gives about 20,000 speakers from a 1993 figure; other sources differ widely. Dialects named are Bero and Zamre. It should not be confused with Masmaje (mes) or the Ethiopian Mesmes (mys). Very little material was found.

\section{Primary Educational Resources}
No primers, textbooks or school programmes in Mesme were found in search results.

\section{Linguistic and Reference Materials}
\begin{itemize}
\item Wikipedia, Mesme language: \url{https://en.wikipedia.org/wiki/Mesme_language}
\item ASJP wordlist entry for Mesme: \url{https://asjp.clld.org/languages/MESME}
\item Numeral-system data sheet for Mesme: \url{https://lingweb.eva.mpg.de/channumerals/Mesme.htm}
\end{itemize}

\section{Religious and Community Resources}
\begin{itemize}
\item Joshua Project, Mesme (Zime) people in Chad: \url{https://www.joshuaproject.net/people_groups/13690/CD}
\end{itemize}

\section{Digital and Audio Resources}
None found in search results.

\section{Curriculum Development Context}
No educational materials were found. Chad's schools use French and Arabic, so a mother-tongue curriculum would need new orthography and materials work with the community.

\end{document}